\documentclass[12pt,a4paper]{article}
\usepackage[margin=2cm]{geometry}
\usepackage[table]{xcolor}
\usepackage[T1]{fontenc}
\usepackage[utf8]{inputenc}
\usepackage{titlesec}
\usepackage{tcolorbox}
\usepackage{enumitem}
\usepackage{tabularx}
\usepackage{fancyhdr}
\usepackage{hyperref}
\usepackage{multicol}
\usepackage{fontawesome5}
\usepackage{graphicx}
\usepackage{tikz}
\usepackage{pifont}
\usepackage{listings}
\usepackage{array}

% ---- COLORS ----
\definecolor{FaranBlue}{HTML}{2563EB}
\definecolor{FaranDark}{HTML}{1E3A5F}
\definecolor{FaranLight}{HTML}{EFF6FF}
\definecolor{FaranGray}{HTML}{6B7280}
\definecolor{FaranAccent}{HTML}{10B981}
\definecolor{FaranYellow}{HTML}{B45309}
\definecolor{CodeBg}{HTML}{1E1E2E}
\definecolor{CodeFg}{HTML}{CDD6F4}
\definecolor{DarkText}{HTML}{111827}
\definecolor{SectionBg}{HTML}{FFFBEB}

\hypersetup{colorlinks=true, linkcolor=FaranYellow, urlcolor=FaranYellow,
    pdftitle={JavaScript Interview Prep - Faran Alam}, pdfauthor={Faran Alam}}

\pagestyle{fancy}
\fancyhf{}
\fancyhead[L]{\small\color{FaranGray}\textbf{Faran Alam} | JavaScript Interview Prep}
\fancyhead[R]{\small\color{FaranGray}\faGlobe\ faran-new-portfolio.vercel.app}
\fancyfoot[C]{\small\color{FaranGray}Page \thepage\ | Faran Alam Portfolio Resources}
\fancyfoot[R]{\small\color{FaranGray}\faWhatsapp\ +92 333 4051830}
\renewcommand{\headrulewidth}{0.4pt}
\renewcommand{\footrulewidth}{0.4pt}
\renewcommand{\headrule}{\hbox to\headwidth{\color{FaranYellow}\leaders\hrule height \headrulewidth\hfill}}
\renewcommand{\footrule}{\hbox to\headwidth{\color{FaranYellow}\leaders\hrule height \footrulewidth\hfill}}

\titleformat{\section}{\Large\bfseries\color{FaranDark}}{\color{FaranYellow}\thesection.}{0.5em}{}
  [\vspace{-4pt}\rule{\textwidth}{1.5pt}\vspace{4pt}]
\titleformat{\subsection}{\large\bfseries\color{FaranYellow}}{\thesubsection.}{0.5em}{}

\tcbuselibrary{skins,breakable}
\newcommand{\circleimage}[3]{%
  \begin{tikzpicture}[baseline=(img.base)]
    \clip (0,0) circle (#1);
    \node[inner sep=0pt] (img) at (0,0) {\includegraphics[width=#2,height=#2]{#3}};
    \draw[line width=0.9pt, color=FaranAccent!85!black] (0,0) circle (#1);
  \end{tikzpicture}%
}
\newtcolorbox{infobox}[1]{colback=SectionBg, colframe=FaranYellow, fonttitle=\bfseries\color{white},
  coltitle=white, colbacktitle=FaranYellow, title=#1, rounded corners, boxrule=1pt, arc=5pt, breakable}
\newtcolorbox{qbox}[1]{colback=FaranLight, colframe=FaranYellow, fonttitle=\bfseries\color{white},
  coltitle=white, colbacktitle=FaranYellow, title={\faQuestionCircle\ Q: #1},
  rounded corners, boxrule=1pt, arc=5pt, breakable}
\newtcolorbox{codebox}[1]{colback=CodeBg, colframe=FaranYellow, fonttitle=\bfseries\color{white},
  coltitle=white, colbacktitle=FaranYellow!80, title={\faCode\ #1}, rounded corners, boxrule=1pt, arc=5pt,
  breakable, left=6pt, right=6pt, top=6pt, bottom=6pt}
\newtcolorbox{tipbox}{colback=green!5, colframe=FaranAccent, fonttitle=\bfseries,
  title={\faLightbulb\ Interview Tip}, rounded corners, boxrule=1pt, arc=5pt, breakable}

\lstset{
  basicstyle=\ttfamily\footnotesize\color{CodeFg},
  backgroundcolor=\color{CodeBg},
  keywordstyle=\color{cyan}\bfseries,
  stringstyle=\color{FaranAccent},
  commentstyle=\color{FaranGray}\itshape,
  breaklines=true, breakatwhitespace=false, columns=fullflexible, keepspaces=true, frame=none, showstringspaces=false
}

\begin{document}

% ==================== TITLE PAGE ====================
\begin{titlepage}
    \pagecolor{FaranDark}\color{white}\centering\vspace*{0.3cm}
    \circleimage{0.8cm}{1.6cm}{logo.jpg}\qquad\circleimage{0.8cm}{1.6cm}{academy-logo.png}\\[0.3cm]
    {\color{FaranYellow}\rule{\textwidth}{3pt}}\vspace{0.8cm}

    \begin{tcolorbox}[colback=FaranYellow, colframe=FaranYellow, width=0.42\textwidth,
        arc=20pt, boxrule=0pt, halign=center]
        \centering\color{white}\textbf{\faBolt\ FREE RESOURCE}
    \end{tcolorbox}
    \vspace{0.8cm}

    {\fontsize{38}{46}\selectfont\textbf{JavaScript}}\\[0.3cm]
    {\fontsize{42}{50}\selectfont\color{FaranYellow}\textbf{Interview Prep}}\\[0.6cm]
    {\large\color{FaranLight}50+ Real Interview Questions with Detailed Answers}\\[0.2cm]
    {\large\color{FaranLight}Core Concepts, Tricky Outputs, \& Coding Challenges}

    \vspace{1.5cm}
    \begin{tcolorbox}[colback=white, colframe=FaranYellow, width=0.68\textwidth,
        arc=8pt, boxrule=1.5pt, halign=center]
        \centering
        \circleimage{1cm}{2cm}{profile.jpg}\\[0.15cm]
        {\large\textbf{\color{FaranDark}Faran Alam}}\\[0.1cm]
        {\color{FaranGray}Full Stack Developer | JavaScript Expert}\\[0.3cm]
        {\small\color{FaranDark}
            \faGlobe\ faran-new-portfolio.vercel.app \quad \faEnvelope\ faran.bsce40@iiu.edu.pk\\[0.1cm]
            \faGithub\ github.com/FaranAlam \quad \faWhatsapp\ +92 333 4051830\\[0.1cm]
            \faFacebook\ facebook.com/faran.alam.5243 \quad \faLinkedin\ linkedin.com/in/faran-alam-14203abc\\[0.1cm]
            \faMapMarkerAlt\ IIU Hostel~6, Islamabad
        }
    \end{tcolorbox}
    \vfill
    {\color{FaranYellow}\rule{\textwidth}{3pt}}
    {\small\color{FaranGray}2026 Edition | Portfolio Resource Series | Faran Alam}
\end{titlepage}
\pagecolor{white}\color{DarkText}
\tableofcontents
\newpage

% ==================== SECTION 1: CORE CONCEPTS ====================
\section{Core JavaScript Concepts}

\begin{qbox}{Explain var, let, and const differences}
\begin{itemize}[leftmargin=1.5em, itemsep=4pt]
    \item \textbf{var}: function-scoped, hoisted to top with value \texttt{undefined}, can be re-declared. \textit{Avoid in modern JS.}
    \item \textbf{let}: block-scoped \texttt{\{\}}, hoisted but NOT initialized (Temporal Dead Zone), re-assignable
    \item \textbf{const}: block-scoped, must be initialized, cannot be re-assigned. Object/array contents can still mutate.
\end{itemize}
\end{qbox}

\begin{qbox}{What is a closure?}
A closure is when an inner function retains access to variables from its outer (enclosing) function's scope, even after the outer function has returned.\\[0.3cm]
\textbf{Classic example -- counter:}
\begin{codebox}{Closure Example}
\begin{lstlisting}
function createCounter() {
  let count = 0;               // outer scope variable
  return {
    increment: () => ++count,  // inner fn remembers count
    decrement: () => --count,
    getCount:  () => count,
  };
}

const counter = createCounter();
counter.increment(); // 1
counter.increment(); // 2
counter.getCount();  // 2  -- count persists!
\end{lstlisting}
\end{codebox}
\textbf{Real use cases:} Private variables, memoization, event handlers, module pattern
\end{qbox}

\begin{qbox}{What is the Event Loop? How does async JS work?}
JavaScript is \textbf{single-threaded}. The Event Loop allows non-blocking operations:
\begin{enumerate}[leftmargin=2em, noitemsep]
    \item \textbf{Call Stack} -- executes synchronous code (top to bottom)
    \item \textbf{Web APIs} -- browser handles async tasks (setTimeout, fetch, DOM events)
    \item \textbf{Microtask Queue} -- resolved Promises (highest priority, runs before macrotasks)
    \item \textbf{Macrotask Queue} -- setTimeout/setInterval callbacks
    \item \textbf{Event Loop} -- moves tasks to stack when it is empty
\end{enumerate}
\begin{codebox}{Event Loop Output Quiz}
\begin{lstlisting}
console.log('1');           // sync - runs first
setTimeout(() => console.log('2'), 0); // macrotask
Promise.resolve().then(() => console.log('3')); // microtask
console.log('4');           // sync

// Output: 1, 4, 3, 2
// Sync first, then microtasks, then macrotasks
\end{lstlisting}
\end{codebox}
\end{qbox}

\begin{qbox}{What is hoisting?}
Hoisting moves declarations (not initializations) to the top of their scope during compile phase.
\begin{codebox}{Hoisting Example}
\begin{lstlisting}
// What you write:
console.log(x);    // undefined (not error!)
var x = 5;

// What JS sees (after hoisting):
var x;             // declaration hoisted
console.log(x);    // undefined
x = 5;             // initialization stays

// let/const: hoisted but NOT initialized
console.log(y);    // ReferenceError: Cannot access 'y'
let y = 10;
\end{lstlisting}
\end{codebox}
\end{qbox}

\newpage

\begin{qbox}{Explain Promises and async/await}
\begin{codebox}{Promises vs Async/Await}
\begin{lstlisting}
// Promise chain (old way)
fetchUser(id)
  .then(user => fetchPosts(user.id))
  .then(posts => console.log(posts))
  .catch(err => console.error(err));

// Async/await (modern way - cleaner)
async function loadUserPosts(id) {
  try {
    const user = await fetchUser(id);
    const posts = await fetchPosts(user.id);
    return posts;
  } catch (err) {
    console.error('Failed:', err.message);
  }
}

// Parallel requests (faster - don't await separately)
const [user, settings] = await Promise.all([
  fetchUser(id),
  fetchSettings(id)
]);
\end{lstlisting}
\end{codebox}
\end{qbox}

\begin{qbox}{this keyword -- what does it refer to?}
\texttt{this} depends on \textit{how} a function is called, not where it is defined.
\begin{itemize}[leftmargin=1.5em, itemsep=4pt]
    \item \textbf{Global scope}: \texttt{window} (browser) / \texttt{undefined} in strict mode
    \item \textbf{Object method}: the object before the dot: \texttt{obj.method()} → \texttt{this = obj}
    \item \textbf{Arrow function}: NO own \texttt{this} -- inherits from enclosing scope
    \item \textbf{Class}: the class instance
    \item \textbf{call/apply/bind}: explicitly set \texttt{this}
\end{itemize}
\end{qbox}

\newpage

% ==================== SECTION 2: ES6+ ====================
\section{ES6+ Features You Must Know}

\begin{center}
\begin{tabularx}{\textwidth}{|>{\bfseries\color{FaranYellow}}l|X|X|}
\hline
\rowcolor{FaranDark}
\color{white}\textbf{Feature} & \color{white}\textbf{What it does} & \color{white}\textbf{Quick Example} \\
\hline
Destructuring & Extract values from arrays/objects & \texttt{const \{name, age\} = user} \\
\hline
\rowcolor{SectionBg}
Spread/Rest & Expand or collect values & \texttt{const copy = [...arr]} \\
\hline
Template Literals & String interpolation & \texttt{`Hello \$\{name\}`} \\
\hline
\rowcolor{SectionBg}
Optional Chaining & Safe property access & \texttt{user?.address?.city} \\
\hline
Nullish Coalescing & Default for null/undefined & \texttt{name ?? 'Anonymous'} \\
\hline
\rowcolor{SectionBg}
Arrow Functions & Short syntax + no \texttt{this} & \texttt{const add = (a,b) => a+b} \\
\hline
Modules & import/export system & \texttt{import \{ fn \} from './utils'} \\
\hline
\rowcolor{SectionBg}
Map / Set & Key-value store / unique values & \texttt{new Map(), new Set()} \\
\hline
Array Methods & Functional array manipulation & \texttt{.map .filter .reduce .flat} \\
\hline
\end{tabularx}
\end{center}

% ==================== SECTION 3: TRICKY OUTPUTS ====================
\section{Tricky Output Questions (Most Asked)}

\begin{infobox}{Predict the Output -- Practice These}
\begin{codebox}{Tricky Question 1: setTimeout in Loop}
\begin{lstlisting}
for (var i = 0; i < 3; i++) {
  setTimeout(() => console.log(i), 0);
}
// Output: 3, 3, 3   (var is function-scoped, loop finishes first)

// FIX using let (block-scoped):
for (let i = 0; i < 3; i++) {
  setTimeout(() => console.log(i), 0);
}
// Output: 0, 1, 2
\end{lstlisting}
\end{codebox}

\begin{codebox}{Tricky Question 2: Object Reference}
\begin{lstlisting}
const a = { x: 1 };
const b = a;           // reference copy, NOT value copy
b.x = 99;
console.log(a.x);      // 99 -- a was also changed!

// To copy: use spread
const c = { ...a };    // shallow copy
c.x = 0;
console.log(a.x);      // 99 -- a is unaffected now
\end{lstlisting}
\end{codebox}

\begin{codebox}{Tricky Question 3: Type Coercion}
\begin{lstlisting}
console.log(1 + '2');    // '12'  (number + string = string)
console.log('5' - 2);   // 3     (string - number = number)
console.log(true + 1);  // 2     (true = 1)
console.log([] + []);   // ''    (empty string)
console.log([] == false); // true  (loose equality!)
// Always use === to avoid coercion bugs
\end{lstlisting}
\end{codebox}
\end{infobox}

\newpage

% ==================== SECTION 4: ARRAY CHALLENGES ====================
\section{Coding Challenges (Write the Solution)}

\begin{codebox}{Challenge 1 -- Reverse a String Without reverse()}
\begin{lstlisting}
function reverseString(str) {
  return str.split('').reduce((acc, char) => char + acc, '');
  // or: return [...str].reverse().join('');
}
reverseString('faran'); // 'naraf'
\end{lstlisting}
\end{codebox}

\begin{codebox}{Challenge 2 -- Remove Duplicates from Array}
\begin{lstlisting}
// Method 1: Set (simplest)
const unique = (arr) => [...new Set(arr)];

// Method 2: filter + indexOf
const unique2 = (arr) => arr.filter((item, idx) => arr.indexOf(item) === idx);

unique([1, 2, 2, 3, 3, 4]); // [1, 2, 3, 4]
\end{lstlisting}
\end{codebox}

\begin{codebox}{Challenge 3 -- Deep Clone an Object}
\begin{lstlisting}
// Simple (no functions/dates):
const deepClone = (obj) => JSON.parse(JSON.stringify(obj));

// With structuredClone (modern browser API):
const clone = structuredClone(obj);

// Recursive (handles edge cases):
function deepCloneCustom(obj) {
  if (obj === null || typeof obj !== 'object') return obj;
  if (Array.isArray(obj)) return obj.map(deepCloneCustom);
  return Object.fromEntries(
    Object.entries(obj).map(([k, v]) => [k, deepCloneCustom(v)])
  );
}
\end{lstlisting}
\end{codebox}

\begin{codebox}{Challenge 4 -- Group Array of Objects by Key}
\begin{lstlisting}
const groupBy = (arr, key) =>
  arr.reduce((groups, item) => {
    const group = item[key];
    groups[group] = groups[group] ?? [];
    groups[group].push(item);
    return groups;
  }, {});

// Usage:
const people = [
  { name: 'Ali', dept: 'Engineering' },
  { name: 'Sara', dept: 'Design' },
  { name: 'Ahmed', dept: 'Engineering' },
];
groupBy(people, 'dept');
// { Engineering: [{Ali...}, {Ahmed...}], Design: [{Sara...}] }
\end{lstlisting}
\end{codebox}

\vspace{0.5cm}
\begin{tipbox}
In interviews, always \textbf{think out loud}. Explain your approach before coding. Interviewers care about \textit{how} you think, not just the answer. If you are stuck, ask: ``Can I clarify the requirements?'' -- this shows professional maturity.
\end{tipbox}


% ==================== ADVANCED BEGINNER + PROFESSIONAL UPGRADE ====================
\section{Beginner-to-Professional Execution Blueprint}

\begin{infobox}{How to Study This Resource the Right Way}
Do not just read this PDF. Convert each major topic into one practical output:
\begin{itemize}[leftmargin=1.5em, itemsep=4pt]
    \item[\ding{52}] Learn the concept for 30--45 minutes
    \item[\ding{52}] Implement one mini feature immediately
    \item[\ding{52}] Write 3--5 bullet notes in your own words
    \item[\ding{52}] Push your progress to GitHub with clear commits
\end{itemize}
This simple loop helps beginners become consistent and interview-ready.
\end{infobox}

\subsection{Weekly Professional Growth Plan}
\begin{center}
\begin{tabularx}{\textwidth}{|>{\bfseries\color{FaranDark}}l|X|>{\centering\arraybackslash}p{2.6cm}|}
\hline
\rowcolor{FaranDark}
\color{white}\textbf{Day} & \color{white}\textbf{Focus} & \color{white}\textbf{Output} \\
\hline
Mon--Tue & Concepts + guided implementation & Notes + code snippets \\
\hline
\rowcolor{SectionBg}
Wed--Thu & Build practical feature/project block & Working module \\
\hline
Fri & Refactor + improve naming/structure & Cleaner codebase \\
\hline
\rowcolor{SectionBg}
Sat & Deploy + test on real device/browser & Live URL + bug list \\
\hline
Sun & Review week + LinkedIn learning post & Public proof of progress \\
\hline
\end{tabularx}
\end{center}

\section{Professional Quality Standards}

\begin{infobox}{Client-Ready Checklist}
Before you call any project complete, verify these essentials:
\begin{itemize}[leftmargin=1.5em, itemsep=4pt]
    \item[\ding{116}] Mobile-first responsive layout
    \item[\ding{116}] Loading, empty, and error states
    \item[\ding{116}] Clean file structure and reusable components
    \item[\ding{116}] Clear README with setup and feature summary
    \item[\ding{116}] Deployed live demo with working links
\end{itemize}
\end{infobox}

\begin{tipbox}
\textbf{Career Shortcut:} Employers do not hire by course certificates. They hire by proof of execution. Ship projects, document your process, and show measurable improvement every month.
\end{tipbox}

\section{30-Day Improvement Sprint}
\begin{itemize}[leftmargin=1.5em, itemsep=5pt]
    \item[\ding{52}] Week 1: Strengthen fundamentals and fix weak concepts
    \item[\ding{52}] Week 2: Build one focused project from scratch
    \item[\ding{52}] Week 3: Add professional polish (SEO, performance, accessibility)
    \item[\ding{52}] Week 4: Deploy, write case study, and publish portfolio update
\end{itemize}
% ==================== LAST PAGE ====================
\newpage
\pagecolor{FaranDark}\color{white}\centering\vspace*{3cm}
{\color{FaranYellow}\rule{\textwidth}{3pt}}\vspace{1cm}
{\fontsize{28}{36}\selectfont\textbf{Practice Daily. Get Hired.}}\\[0.5cm]
{\large Explore more resources at \textcolor{FaranYellow}{faran-new-portfolio.vercel.app}}
\vspace{0.4cm}
\circleimage{1.25cm}{2.5cm}{academy-logo.png}\\[0.4cm]
\begin{tcolorbox}[colback=FaranYellow, colframe=FaranYellow, width=0.72\textwidth,
    arc=10pt, boxrule=0pt, halign=center]
    \centering\color{white}
    {\large\textbf{Faran Alam}}\\[0.3cm]
    \faGlobe\ faran-new-portfolio.vercel.app \quad \faEnvelope\ faran.bsce40@iiu.edu.pk\\[0.2cm]
    \faWhatsapp\ +92 333 4051830 \quad \faGithub\ github.com/FaranAlam\\[0.2cm]
    \faFacebook\ facebook.com/faran.alam.5243 \quad \faLinkedin\ linkedin.com/in/faran-alam-14203abc\\[0.2cm]
    \faMapMarkerAlt\ IIU Hostel~6, Islamabad
\end{tcolorbox}
\vspace{1.5cm}
{\color{FaranYellow}\rule{\textwidth}{3pt}}
{\small\color{FaranGray}Portfolio Resource Series | Faran Alam | 2026}

\end{document}
